% --------------------------------------------------------------------------
% File: src/architecture/decisions/templates/adr-software-structure-template.tex
% Style: style v1.0 (tooling/latex/style.sty)
% Engine: pdfLaTeX through tooling/scripts/latex_build.py
% Purpose: reusable unfilled template; not an approved project decision
% --------------------------------------------------------------------------
\documentclass[11pt,letterpaper]{article}
\usepackage{technical-design-spec}
\renewcommand{\DocTitle}{ADR: Software Structure and Component Boundaries}
\renewcommand{\DocSubtitle}{Modularity, ownership, dependency direction and changeability}
\renewcommand{\DocVersion}{v1.0}
\renewcommand{\DocStatus}{Template}
\renewcommand{\DocOwner}{Architecture Governance}
\renewcommand{\DocAuthor}{Architecture Documentation Library}
\renewcommand{\DocDate}{2026-10-28}
\renewcommand{\DocSummary}{An unfilled, domain-specific architecture decision record for software structure and component boundaries. Inherits the central decision structure and adds specialized design evidence.}
\begin{document}
\makedoctocpage


\section{Record identity and decision authority}
\noindent\textbf{Status:} \ADRStatus{Proposed}\quad
\textbf{ADR ID:} \emph{[ADR-0000]}\quad
\textbf{Date:} \emph{[YYYY-MM-DD]}

\begin{description}
\item[Decision title:] \emph{[One sentence describing a single architecturally significant decision.]}
\item[Entity of interest:] \emph{[System, platform, business capability, product or bounded context; stable identifier.]}
\item[Scope and baseline:] \emph{[Repository, service, environment(s), architecture-description ID and revision.]}
\item[Owner and accountable approver:] \emph{[Roles or teams; decision authority, approval evidence and date.]}
\item[Consulted stakeholders:] \emph{[Security, operations, users, product, architecture, compliance and others as applicable.]}
\item[Related records:] \emph{[Links to dependent/conflicting decisions; supersedes / superseded by, if applicable.]}
\item[Review and expiry trigger:] \emph{[Change in assumptions, threshold or policy, or next review milestone.]}
\end{description}

\begin{notebox}[Template, not an approval]
Replace all bracketed guidance before requesting review. The \texttt{Proposed}
badge is illustrative. A decision is not \texttt{Accepted} without recorded
authority and evidence. Preserve historical accepted records: supersede with
a new ID instead of silently rewriting their original rationale.
\end{notebox}

\section{Context and architectural significance}
\begin{description}
\item[Problem and forces:] \emph{[What is changing, why now, and why this cannot be left to local implementation discretion?]}
\item[Stakeholders and concerns:] \emph{[Stable stakeholder / concern IDs and impacted quality attributes; reference architecture views.]}
\item[Drivers and constraints:] \emph{[Functional and nonfunctional requirements, security policies, standards, contracts, cost and schedule limits.]}
\item[Boundary and assumptions:] \emph{[Included and excluded responsibilities, dependencies and assumptions with owners.]}
\item[Decision deadline:] \emph{[When must the decision take effect; what happens if deferred?]}
\end{description}

\section{Options considered and evidence}
Include the status quo as an explicit alternative where practical. Evaluate
at least two feasible options, or explain why only one remains viable.
\begin{tabularx}{\linewidth}{@{}p{0.18\linewidth}X X@{}}
\toprule
\textbf{Option} & \textbf{Benefits and supporting evidence} & \textbf{Costs, risks and reason for disposition}\\
\midrule
A: \emph{[baseline]} & \emph{[measured outcome; evidence URI]} & \emph{[trade-off; accept/reject]}\\
B: \emph{[candidate]} & \emph{[measured outcome; evidence URI]} & \emph{[trade-off; accept/reject]}\\
C: \emph{[optional]} & \emph{[measured outcome; evidence URI]} & \emph{[trade-off; accept/reject]}\\
\bottomrule
\end{tabularx}

\noindent\textbf{Evaluation method:} \emph{[Decision criteria, weights if used,
disqualifiers, assumptions, sensitivity analysis and evidence quality. Avoid
presenting subjective scoring as objective fact.]}

\section{Decision, rationale and non-goals}
\begin{description}
\item[We choose:] \emph{[Precise decision, including scope, interface boundaries and effective version.]}
\item[Because:] \emph{[Why the chosen option best satisfies the stated drivers; cite trade-off evidence.]}
\item[Non-goals:] \emph{[Capabilities intentionally not decided or implemented here.]}
\item[Invariants and exceptions:] \emph{[Conditions that must remain true; exceptional cases and their owners.]}
\item[Dependencies:] \emph{[Required predecessor ADRs, roadmap items, provider dependencies and migration prerequisites.]}
\end{description}

\section{Domain-specific decision inputs}
Complete the following prompts with referenceable designs, measurements,
threat assessments or operational constraints; mark inapplicable prompts
explicitly rather than removing necessary decision context.
\begin{itemize}
\item Domain boundaries, components, team ownership, release cadence and data ownership.
\item Modular monolith versus services; layered versus event-driven versus plugin; runtime versus build-time composition.
\item Interface stability, dependency direction, cyclic imports, shared code/versioning and distributed transactions.
\item Changeability, test seams, runtime isolation, deployment autonomy and security boundary enforcement.
\item Migration plan: strangler increments, compatibility layer, ownership transfer and decommissioning.
\end{itemize}
\section{Domain-specific alternatives and trade-offs}
Explain how each feasible option addresses the following criteria.
\begin{itemize}
\item Cognitive load and coupling, delivery bottlenecks, deployment independence and runtime complexity.
\item Network and data consistency penalties versus modularity and independent scale.
\item Proof of boundary enforcement through dependency analysis and architecture tests.
\end{itemize}
\noindent\textbf{Domain comparison:} \emph{[Attach a scored/weighted matrix only
when scoring assumptions are explicit; link to proofs and rejected options.]}
\section{Domain-specific validation and review gates}
The proposed decision is ready for acceptance only after its applicable
tests have objective pass/fail thresholds, an owner and dated evidence.
\begin{itemize}
\item Dependency rules enforced in CI; representative feature change confined to intended module.
\item Integration/contract tests under compatible release combinations; performance comparison.
\item Migration and rollback demonstration, team ownership confirmation and architectural fitness checks.
\end{itemize}
\noindent\textbf{Related repository material:} Views and Beyond architecture styles, design patterns and micro-frontend evaluation.

\section{Consequences, risks and safeguards}
\begin{description}
\item[Positive consequences:] \emph{[Expected measurable benefit; beneficiary and timeframe.]}
\item[Negative consequences:] \emph{[Accepted technical debt, cost, complexity, coupling and operational burden.]}
\item[Security and privacy effects:] \emph{[Threat exposure, data protection, identity, trust boundaries and residual risk owner.]}
\item[Operational and financial effects:] \emph{[Support ownership, on-call impact, capacity, license or cloud spending.]}
\item[Unknowns and risk treatment:] \emph{[Risk IDs, mitigation, explicit acceptance authority, or experiment planned.]}
\end{description}

\section{Implementation, verification and reversibility}
\begin{tabularx}{\linewidth}{@{}p{0.22\linewidth}X@{}}
\toprule
\textbf{Control} & \textbf{Complete with concrete, reviewable evidence}\\
\midrule
Implementation owner and plan & \emph{[Work-item IDs, milestones, responsible team and migration/rollout path.]}\\
Acceptance tests & \emph{[Functional and quality-attribute test IDs, measurable thresholds and environments.]}\\
Security / reliability gates & \emph{[Control IDs, negative tests, sign-off conditions and exception process.]}\\
Evidence and provenance & \emph{[Artifact URI, revision/commit, test date, result and reviewer.]}\\
Rollback or exit & \emph{[Trigger, data/recovery considerations, owner, estimated reversal window.]}\\
Post-deployment review & \emph{[Success metric, observation interval, reviewer, revisit deadline.]}\\
\bottomrule
\end{tabularx}

\section{Architecture-description traceability}
A decision record contributes rationale to an architecture description; it
does not, by itself, constitute a conforming ISO/IEC/IEEE 42010:2022
architecture description.
\begin{description}
\item[Description and configuration baseline:] \emph{[AD identifier, repository URI, revision.]}
\item[Stakeholders and concerns:] \emph{[SH-*, CN-* and their source register.]}
\item[Viewpoints, views and model kinds:] \emph{[VP-*, VW-*, VC-* and governing specifications.]}
\item[Correspondence and consistency:] \emph{[CO-*, CM-*; checks, known inconsistency and disposition.]}
\item[Requirements and evidence:] \emph{[Requirement/control IDs, test or audit artifacts and link revisions.]}
\end{description}

\section{Review history and references}
\begin{tabularx}{\linewidth}{@{}p{0.16\linewidth}p{0.19\linewidth}p{0.19\linewidth}X@{}}
\toprule
\textbf{Date} & \textbf{Reviewer / role} & \textbf{Disposition} & \textbf{Evidence and requested changes}\\
\midrule
\emph{[date]} & \emph{[role]} & \emph{[pending]} & \emph{[review or PR URL]}\\
\bottomrule
\end{tabularx}

\noindent\textbf{Sources:} \emph{[Link to relevant requirements, architecture
models, policies, benchmark results, contracts, standards and dated evidence.]}

\end{document}
